% Generated by Listas/make_statements.py from pratica-lecture-06.md, item 6. Do not edit.
\textbf{Source.} Milgrom \& Roberts (1992), \emph{Economics, Organization and Management}, Ch. 8 Exercises, Food for Thought 3, p.~284 {[}PDF 299{]}. Reproduced verbatim. Answer in the source: none in the book.



\begin{quote}

\begin{enumerate}

\def\labelenumi{\arabic{enumi}.}

\setcounter{enumi}{2}

\tightlist

\item

  In many countries, corporations whose shares are publicly traded are required to have their financial statements audited by independent accountants who check whether the financial information being provided by management to investors is accurate and has been prepared following accepted methods and procedures and who then publicly attest to their findings. The auditors are generally chosen by management (perhaps subject to nominal stockholder approval). Audit work provides a major source of income for accounting firms. Generally, each accounting firm has many clients that it audits: Even when it might be technically feasible for a relatively small accounting firm to audit a large corporation's records, there seems to be a reluctance for corporations to use audit firms when a single client would represent too much of the accountant's business. Instead, a handful of extremely large accounting firms typically do almost all the auditing of large corporations, with each having many corporate clients. How do you account for this?

\end{enumerate}

\end{quote}

